\documentclass[11pt]{article}
\usepackage[margin=0.88in]{geometry}
\usepackage{booktabs,tabularx,array,longtable,microtype,xcolor}
\usepackage{amsmath,amssymb}
\usepackage[hidelinks]{hyperref}
\usepackage[numbers,sort&compress]{natbib}
\usepackage{enumitem}
\setlist{nosep,leftmargin=*}
\newcommand{\sourcearm}{\textsc{Source}}
\newcommand{\broadcastarm}{\textsc{Commander--Broadcast}}
\newcommand{\stararm}{\textsc{Commander--Star}}
\newcommand{\scp}{\texttt{SCP1}}
\newcommand{\code}[1]{\texttt{#1}}
\title{An Embodied Squadron Commander for Long-Horizon LLM Teams\\
Related Work, Protocol Design, and Auditable Evaluation Plan}
\author{AlemDICE Research Engineering Report}
\date{22 July 2026}

\begin{document}
\maketitle

\begin{abstract}
Language-model agents can interpret new objectives and improvise in unfamiliar
worlds, but unconstrained peer discussion does not reliably produce one coherent
team plan.  This report develops a deliberately narrow alternative: one embodied
agent is the squad commander and owns strategic task planning, while two wingmen
retain tactical autonomy for movement, prerequisites, immediate survival, and
execution.  The commander is not a privileged fourth process and receives no
global simulator state.  A separate serial model call uses only the commander's
legal text observation and validated prior reports to issue a versioned,
short-leased plan; the same three embodied agents then choose actions in parallel.
The design is inspired by joint-intention and SharedPlans theories, practical
teamwork models, role-based multi-agent learning, LLM supervisor--worker systems,
embodied language planning, and distributed leases.  It differs from all of
these in an important way: the scientific treatment is an online authority and
communication protocol layered onto an otherwise unchanged Alem Source policy.
We specify failure semantics, information boundaries, metrics, preregistered
30-step and three-seed 200-step Luna studies, and a path to multiple squads.  No
new hierarchy efficacy result is claimed in this report; hosted-model outcomes
remain pending until the staged gates and semantic trace audit complete.
\end{abstract}

\section{Research question and thesis}

The motivating question is not merely whether three LLM agents can exchange
useful text.  It is whether a team can maintain a coherent division of labor over
many embodied decisions without silently changing the environment, action policy,
or model underneath the comparison.  Alem provides persistent, partially
observable multi-agent interaction, differentiated physical specializations,
networked language-model policies, exact trajectories, and achievement-level
outcomes~\citep{alem2026}.  It therefore exposes a recurring gap between a plan
that sounds coordinated and behavior that actually completes joint work.

The primary thesis is:

\begin{quote}
For one small squad, separating strategic authority from tactical execution can
improve plan coherence and role persistence, provided that authority is explicit,
plans expire, reports are validated, local action choice remains grounded in legal
observations, and the serial planning overhead is measured rather than ignored.
\end{quote}

``Squadron'' is an organizational metaphor, not a claim that military doctrine is
being reproduced.  The operative abstraction is one commander and two wingmen in
a single team.  The commander has a body, fixed Alem specialization, observations,
actions, rewards, and risks like every other player.  It makes two logically
separate calls when planning is due: first a planning-only call, then its ordinary
embodied action call in parallel with the other two agents.  The planner cannot
act; the action call cannot revise the plan.  Future work may compose several
squads, but no cross-squad mechanism is part of the present claim.

Three observations motivate this direction.  First, previous local peer-protocol
pilots produced positive cohesion signals on isolated trajectories, but consensus
often failed to become synchronized execution and prompt-based role auctions were
fragile.  Second, an earlier bodyless-leader treatment added a fourth logical
participant and exposed every worker view to it; that is a useful comparison, but
it is not the proposed embodied command structure.  Third, the public Source
behavior already has valuable properties---three concurrent action calls,
affordance-grounded action parsing, private scratchpads, and one-tick peer
broadcast---which should remain intact in the control arm.

\section{Related work}

\subsection{Joint commitments and explicit team plans}

Joint-intention theory treats teamwork as more than coincident individual goals:
participants maintain commitments to a joint activity and have obligations when
its status changes~\citep{cohen1991teamwork}.  SharedPlans develops a richer account
of partial plans, individual subplans, mutual belief, and the relationships among
contributors to a complex group action~\citep{grosz1996sharedplans}.  These theories
suggest why an LLM's free-form ``I will help'' is too weak for measurement.  It
does not identify a plan version, task owner, dependency, success condition, or
the point at which the statement ceases to authorize action.

STEAM translated teamwork theory into a reusable agent architecture for dynamic,
uncertain domains.  Its central contribution was not a particular domain plan but
domain-independent rules for establishing and maintaining team activity under
partial and inconsistent views~\citep{tambe1997steam}.  Later decision-theoretic
analysis formalized the trade-off between communication cost and coordination
benefit~\citep{nair2005team}.  The proposed squad protocol follows this tradition
by making commitments and termination observable.  It is intentionally less
ambitious epistemically: it does not claim mutual belief or a formal team-optimal
policy.  A parsed record is evidence that a message satisfied a protocol, not that
the model truly understood or obeyed it.

Contract Net is another relevant ancestor.  It distributes tasks through
announcement, bidding, award, and contractor execution rather than assuming a
single fixed allocator~\citep{smith1980contractnet}.  That mechanism remains
important for future dynamic team formation, but it did not work reliably in the
initial prompt-only Alem role pilot: self-bids were poorly calibrated, awards
could violate the visible ranking, and acknowledgements appeared without valid
awards.  The present method therefore tests fixed, explicit command authority
before returning to negotiated formation.  It narrows the causal question from
``can every mechanism coexist?'' to ``does one valid planner plus bounded reports
improve an otherwise matched team?''

\subsection{Organization, hierarchy, and command failure}

Multi-agent organizational research distinguishes hierarchies, teams, markets,
coalitions, congregations, and other structures by authority, information flow,
adaptability, and failure modes~\citep{horling2004organizations}.  A hierarchy can
reduce conflicting decisions and communication fan-out, but creates a bottleneck
and a concentrated failure point.  A team model can support richer mutual
adjustment, but may spend more computation establishing common ground.  This is
why broadcast and star routing are separated from the authority treatment: the
same commander plan should be testable with or without peer-visible statuses.

The design borrows the \emph{expiration} idea from distributed leases.  Gray and
Cheriton used time-bounded rights so failures affect availability rather than
allowing stale authority to persist indefinitely~\citep{gray1989leases}.  A squad
task lease is not a distributed cache-consistency proof, and simulator ticks avoid
the clock-skew assumptions of physical-time leases.  The analogy nevertheless
provides a clean safety rule: an invalid or missing review cannot silently renew
old authority.  If repeated invalid plans exhaust a lease, wingmen stop pursuing
the expired task and report that they are waiting.  Version and team identifiers
make delayed or cross-team records rejectable.

The first implementation deliberately has no failover.  Promoting a wingman after
commander failure would alter authority, introduce an election mechanism, and make
the initial result harder to interpret.  Commander inactivity, planning transport
failure, invalid output, and plan expiry are logged separately.  A provider
transport failure retains the existing evaluator's fail-and-resume semantics; a
semantic planning failure preserves only the unexpired last valid plan.

\subsection{Centralized training and hierarchical multi-agent learning}

Centralized training with decentralized execution (CTDE) addresses learning
non-stationarity by exposing a critic or value mixer to joint information during
training while keeping deployed actors local.  MADDPG learns centralized critics
with per-agent actors~\citep{lowe2017maddpg}; QMIX factorizes a joint value subject
to a monotonicity constraint that permits decentralized greedy action
selection~\citep{rashid2018qmix}.  These are not direct instances of the proposed
method.  The squad commander performs \emph{online} inference at execution time,
changes the prompts presented to actors, and incurs serial latency and tokens.
Calling it CTDE would hide precisely the intervention being measured.

Role-oriented MARL asks how roles can reduce a joint problem.  ROMA learns
identifiable emergent role representations~\citep{wang2020roma}; RODE constructs a
bi-level hierarchy in which a lower-frequency selector chooses restricted role
action spaces and role policies operate at the primitive level~\citep{wang2021rode}.
This temporal decomposition strongly motivates a five-tick planning cadence, but
there are two crucial differences.  Alem roles such as warrior, forager, and miner
are fixed environment capabilities, whereas squad assignments are temporary
responsibilities.  Moreover, the LLM planner emits language tasks rather than a
learned latent role or restricted discrete action set.  Compliance must therefore
be measured from messages and behavior rather than guaranteed by policy
architecture.

\subsection{Supervisor--worker LLM systems and organizational topology}

LLM multi-agent work shows several ways to impose structure.  CAMEL uses
role-playing and inception prompts to sustain cooperative conversations
~\citep{li2023camel}.  AutoGen exposes programmable conversational agents and
interaction patterns~\citep{wu2023autogen}.  MetaGPT encodes standardized workflows
and assigns specialist roles in a software-production pipeline
~\citep{hong2024metagpt}.  DyLAN dynamically selects agents and communication links
for each task~\citep{liu2023dylan}.  Collectively, these systems show that topology
and role prompts are part of the algorithm, not incidental user-interface choices.

Most such evaluations are disembodied question answering, coding, or artifact
production.  A supervisor can often inspect a shared transcript or a worker's
finished artifact; an embodied commander instead has a partial local view, workers
act simultaneously, and a plausible instruction may be physically impossible at
that tick.  The Alem method also requires a strict control: the baseline must not
receive supervisor vocabulary, structured status demands, extra memory, or
different sampling parameters.  These constraints turn an orchestration pattern
into a falsifiable behavioral treatment.

\subsection{Embodied language planning and multi-robot coordination}

Language-model planning for embodied agents reveals why strategy and affordance
grounding should remain separate.  Zero-shot planners can extract useful action
knowledge but require admissibility filtering and replanning in the environment
~\citep{huang2022zeroshot}.  SayCan combines language-based task relevance with
value-derived skill feasibility, preventing semantically appealing yet infeasible
steps from dominating execution~\citep{ahn2023saycan}.  ReAct interleaves reasoning
and action so environmental feedback can revise a plan~\citep{yao2023react}.  In
Alem, the Source affordance list and action validator continue to ground each
physical action; the squad planner never bypasses them.

SMART-LLM decomposes instructions, forms coalitions, and allocates work for
multiple robots~\citep{kannan2023smartllm}.  RoCo lets robot-specific LLM agents
discuss strategies and combines their high-level output with motion planning and
environment feedback~\citep{mandi2023roco}.  Both support the premise that language
models can help with multi-robot task structure.  Our experiment asks a different
question: under matched repeated interaction and a fixed team, is a single local
commander more reliable and economical than leaving strategic planning to all
three agents?  The answer cannot be inferred from one-shot plan quality.

\subsection{Memory, context, and long-horizon role coherence}

Long missions create two distinct memory problems.  An agent must remember world
facts that have left its observation window, and it must retain the currently
authorized responsibility without letting old instructions compete with new ones.
Generative Agents uses an experience stream, retrieval, reflection, and planning
to maintain behavior over time~\citep{park2023generative}.  Reflexion stores
linguistic feedback in episodic memory to improve later decisions
~\citep{shinn2023reflexion}.  MemGPT treats the context window as a managed memory
hierarchy with explicit movement between tiers~\citep{packer2023memgpt}.

The first squad treatment does not add a retrieval database.  It keeps Source's
private worker scratchpad unchanged, adds a bounded planner-only scratchpad, and
places a compact authoritative plan in every current action prompt.  Plan version,
lease, and assignment are middleware state rather than memories the model must
reconstruct.  This distinction matters: a \emph{memory hierarchy} manages which
information is available, while a \emph{command hierarchy} determines who may
change team objectives.  Future experiments can cross these factors instead of
assuming that better memory automatically produces role compliance.

\subsection{Resulting gap}

The literature establishes strong components but not the full target combination:
one embodied, locally informed commander; unchanged concurrent action policies;
atomic structured plans; explicit authority; bounded leases; validated delayed
reports; and paired long-horizon evaluation against a source-matched control.
Table~\ref{tab:distinctions} states the main boundaries.

\begin{table}[t]
\centering\small
\begin{tabularx}{\textwidth}{lXXX}
\toprule
Area & Central mechanism & What transfers & What must not be conflated\\
\midrule
Joint plans & commitments and subplans & explicit responsibility and termination & parsed text with mutual belief\\
CTDE & global training signal, local actors & compare centralized information assumptions & online planner with training-only critic\\
Hierarchical MARL & slow role selector, fast policies & temporal abstraction & latent role with language assignment\\
LLM workflows & supervisor and specialist agents & authority and topology as algorithm & shared artifact pipeline with embodied control\\
Memory systems & select, summarize, retrieve & bounded context state & memory tier with command authority\\
Distributed leases & expiring rights & stale-authority containment & simulator protocol with cache-consistency proof\\
\bottomrule
\end{tabularx}
\caption{Conceptual distinctions needed for a precise research claim.}
\label{tab:distinctions}
\end{table}

\section{Method: one embodied squad}

\subsection{Participants and information boundary}

Let the physical team be $M=\{0,1,2\}$ and let $c\in M$ be configurable.  The
initial study fixes $c=0$.  There are exactly three action clients and three
trajectory slots.  When planning is due, a separate client instance is cloned
from client $c$ with the same provider, model, generation settings, timeout,
retries, and context policy, but a planner-specific prompt-cache suffix.  Its
usage is attributed to agent $c$ under a distinct planning phase.

At tick $t$, the planner input is limited to:
\begin{enumerate}
  \item the public game rules;
  \item agent $c$'s legal long- and short-term text observation at $t$;
  \item the last accepted plan and its metadata;
  \item a bounded private planner scratchpad; and
  \item previously validated \scp{} reports from members.
\end{enumerate}
It does not receive other agents' direct observations, raw JAX state, reward
internals, analysis-only opportunity labels, worker scratchpads, or worker
reasoning.  This is a local commander, not an omniscient coordinator.

\subsection{Plan schema and state machine}

A planning response contains one tagged JSON object.  \code{REPLACE} requires a
nonempty objective and exactly one assignment for every member.  An assignment
contains \code{agent\_id}, unique \code{task\_id}, \code{directive},
\code{target}, observable \code{completion}, dependency task IDs, and an optional
synchronization action/window.  Dependencies must resolve within the plan and
form an acyclic graph.  Synchronization actions must be canonical Alem actions
and their windows must fit inside the lease.  Validation is atomic: one malformed
assignment rejects the proposal.

Middleware, not the model, adds team ID, monotonic version, issue tick, next
review tick, and expiry tick.  A scheduled review may return \code{KEEP} only
when a plan is active; this preserves the version and renews the lease.  Event
reviews and invalid-plan retries require \code{REPLACE}.  An invalid response
does not change or extend the prior plan.  If no unexpired plan remains, agents
may take an immediately necessary survival action; otherwise they choose
\code{Noop} and report \code{WAITING}.

Planning occurs at tick zero, every five ticks, on the next tick after a new
\code{BLOCKED}, \code{COMPLETE}, or \code{EMERGENCY} transition, and on the next
tick after invalid plan output.  Repeated identical status states are deduplicated.
Unscheduled replans are capped at $\lceil T/10\rceil$, yielding at most nine plan
calls in a 30-step episode and 60 in a 200-step episode.  A scheduled call may
absorb a pending event without adding an extra call, but it remains an event
review and cannot use \code{KEEP}.

\subsection{Execution authority}

After validation, the full accepted plan is placed in all three action prompts
before the same tick's parallel action calls; each prompt highlights its owner's
assignment.  This shared plan view prevents dependencies from becoming opaque.
Only the planning phase can change objectives or assignments.  The commander's
action phase executes its own task and cannot revise the plan.  Wingmen may choose
local routes, acquire prerequisites, react to immediate danger, and determine
feasible primitive actions, but may not create competing objectives, reassign
members, or treat scratchpad plans as team authority.

This is a prompt-level behavioral contract, not forced action substitution.  The
existing action generator, visible reasoning behavior, scratchpad, affordance
display, semantic-retry policy, parser, and environment validity check remain in
the path.  Consequently, failure to follow an assignment remains observable and
scientifically meaningful rather than being hidden by middleware-selected actions.

\subsection{Status records and routing}

Agents use the existing optional \code{<communication>} field for one bounded
record:
\begin{center}
\small\code{SCP1|TEAM=squad-0|VERSION=3|TASK=ore-2|STATE=ACTIVE|EVIDENCE=at iron vein|REQUEST=wood}
\end{center}
Allowed states are \code{ACK}, \code{ACTIVE}, \code{BLOCKED}, \code{COMPLETE},
\code{EMERGENCY}, and \code{WAITING}.  The ledger rejects malformed records,
wrong team IDs, stale versions, and tasks not assigned to the sender.  Rejected
records remain in the raw audit trail but cannot trigger planning.  A worker's
attempt to embed a plan is counted and never accepted.

\broadcastarm{} retains Source's one-tick peer broadcast and also enters valid
statuses in the commander ledger.  \stararm{} gives the same planner identical
reports and plan timing, but does not deliver wingman statuses to peers.  This
isolates routing from strategic authority.  \sourcearm{} retains ordinary free
broadcast and has no commander instructions, schema, or planner call.

\subsection{Failure semantics}

\begin{longtable}{p{0.23\textwidth}p{0.31\textwidth}p{0.36\textwidth}}
\toprule
Failure & Runtime response & Scientific record\\
\midrule
Provider transport error & fail episode; preserve attempt; resume seed later & attempts, error types, requests, billed usage\\
Invalid plan JSON/schema & retain only unexpired plan; retry next tick within cap & raw output, rejection, unchanged expiry\\
Invalid \code{KEEP} & reject; require replacement & invalid-keep counter\\
Expired plan & remove authority; survival/Noop and \code{WAITING} & uncovered tick and expiry counter\\
Malformed/stale status & audit but do not enter active ledger or trigger review & reason-specific rejection counter\\
Commander inactive & planner still uses its legal inactive text; action constraints remain & inactivity and plan behavior visible\\
Commander unavailable & no failover in v1 & explicit single-point failure\\
Wingman noncompliance & do not override primitive action & assignment/action trace for semantic review\\
\bottomrule
\end{longtable}

\section{Implementation plan and delivered architecture}

The implementation is additive.  Existing \code{baseline}, \code{leader\_peer},
and \code{leader\_no\_peer} behavior is preserved; the new topologies are
\code{embodied\_commander\_broadcast} and
\code{embodied\_commander\_star}.  The one-team configuration is already
namespaced by \code{team.id} so later squads do not require rewriting the record
format.

\begin{table}[h]
\centering\small
\begin{tabularx}{\textwidth}{lX}
\toprule
Component & Responsibility\\
\midrule
\code{team\_commander.py} & strict plan/status parsing, dependency validation, leases, triggers, authority ledger, planner prompt, directives, metrics\\
\code{AgentFactory} & clone commander client without adding a fourth client; isolate cache namespace\\
\code{evaluator.py} & serial due-plan phase, same-tick directive delivery, three parallel action calls, routing and durable accounting\\
\code{robust\_all.py} & treatment-only authority/status instructions; Source prompt unchanged\\
Hydra profiles & exact Luna Source settings with dedicated 30/200 seeds and caps\\
Study launcher & preflight, dry run, immutable manifest, hashes, call ceilings, resume, staged arms\\
Study summarizer & paired episode table, attempt-ledger usage, explicit availability, preregistered gates\\
Tests & parser/state properties, routing identities, leak sentinels, parallel ordering, trajectory parity, profile and cache checks\\
\bottomrule
\end{tabularx}
\caption{Implementation boundaries.  File names are relative to the AlemDICE LLM baseline.}
\end{table}

The following invariants are release blockers:
\begin{enumerate}
  \item \textbf{Control preservation.} A scripted \sourcearm{} episode and a
  scripted commander episode with identical primitive actions have identical
  observations, actions, rewards, dones, and physical trajectory dimensions.
  \item \textbf{No fourth participant.} Commander treatments report three
  physical and three logical participants; planning usage belongs to the selected
  embodied commander.
  \item \textbf{Same-tick order.} An accepted plan is visible before all three
  action calls, and a concurrency barrier confirms those calls overlap.
  \item \textbf{No hidden-state path.} Sentinel raw-state fields and non-commander
  legal views never appear in the planner prompt.
  \item \textbf{Atomic authority.} Invalid plans, stale statuses, wrong tasks, and
  worker-authored plans cannot modify the active ledger or extend a lease.
  \item \textbf{Auditability.} Resolved base/per-arm configuration hashes, prompt
  contract hash, source and lock hashes, cache keys, exact commands, attempts,
  failures, and usage are durable.
\end{enumerate}

\section{Evaluation design}

\subsection{Arms and held-constant variables}

The initial paid experiment contains two arms only:
\begin{description}
  \item[\sourcearm{}:] current Source protocol: three \code{robust\_all}
  agents, ordinary peer broadcast, private scratchpads, and concurrent actions.
  \item[\broadcastarm{}:] the full structured treatment: Agent 0 planner,
  leased plan, executor authority, validated \scp{} statuses, and retained peer
  broadcast.
\end{description}
The model is \code{gpt-5.6-luna} through the OpenAI Responses API, with
temperature 1.0, reasoning effort \code{none}, 8192 output tokens, disabled
storage, three prompt-cache traffic shards, explicit 30-minute cache mode,
420-second timeout, five transport attempts, and identical worker role prompts.
The environment, difficulty, seed, action code, context lengths, rendering,
retry policy, and debrief setting are matched.  Cache namespaces differ by arm
to prevent a condition from accidentally inheriting another condition's cached
prefix accounting.

\stararm{} is deferred until broadcast passes.  Commander IDs 1 and 2 are then
tested on the winning routing to determine whether results depend on Agent 0's
warrior specialization or map position.  This ordering avoids spending the
initial budget on a three-way interaction before the core authority mechanism is
known to work.

\subsection{Stage 0: deterministic and provider preflight}

Static and scripted checks precede paid calls.  The full test suite and Ruff must
pass; the two new profiles must compose; launcher dry runs must create no files;
and the deterministic smoke episode must preserve the Source trajectory contract.
The hosted preflight makes exactly two bounded Luna requests: one response that
must contain a parseable \code{Noop} action and one response that must contain a
valid three-member \code{REPLACE} plan.  A preflight is compatibility evidence,
not behavioral evidence.

\subsection{Stage 1: 30-step wiring and qualitative gate}

Both arms run Easy seed 12000 for 30 ticks.  Promotion requires:
\begin{itemize}
  \item both canonical episodes complete with zero transport errors;
  \item at least one valid treatment plan;
  \item active plan coverage $\geq 0.80$;
  \item action parse rate $\geq 0.95$ and attempted-status parse rate $\geq 0.90$;
  \item zero accepted unauthorized plans, zero accepted stale assignments, and
  zero construction-level hidden-state guard violations; and
  \item manual review of planner and action traces, including one invalid or
  transition case if present.
\end{itemize}
Passing this stage means the mechanism is executable and auditable.  A return or
achievement difference at one short seed is not an efficacy result.

\subsection{Stage 2: three paired seeds at 200 steps}

Before inspection, both arms are frozen on Easy seeds 12000--12002 and 200 ticks.
The exploratory promotion rule requires all three pairs valid, achievement rate
higher on at least two seeds, an action-parse decline larger than two percentage
points on at most one seed, at least 90\% plan coverage, at least 90\% valid plan
calls and attempted status reports, zero accepted authority/staleness violations,
and no more than 35\% mean regression in either total tokens or wall time per
step.  All seed-level values are reported even if the aggregate gate fails.

Three seeds are too few for a confirmatory claim.  This stage decides whether the
mechanism merits the next ablations, diagnoses overhead, and estimates variance.
If it fails, revisions move to new development seeds and receive a new protocol
version; the original outputs remain immutable.

\subsection{Stage 3 and confirmatory expansion}

If broadcast passes, run star routing on the same development seeds, then rotate
the commander identity.  Freeze the winning combination and evaluate at least 20
held-out paired seeds per Easy, Medium, and Hard condition.  Power calculations
should use the observed paired variance, not an arbitrary universal seed count.
Report paired bootstrap intervals for continuous outcomes, exact paired results
for sparse binary achievements where appropriate, and a hierarchical model with
seed and difficulty effects.  No across-difficulty average replaces the three
difficulty-specific estimates.

Long-horizon evaluation crosses 200, 500, 1000, and environment-terminal caps.
It should measure role consistency inside opportunity windows, stale-task
execution after replacement, dependency satisfaction, synchronized-action
success, blocker recovery, plan churn, and assignment switches.  A later
factorial experiment crosses command structure with typed event memory, because
authority and memory are separate hypotheses.

\subsection{Metrics and causal diagnostics}

Task outcomes include team return, total/normal/coordination achievement rates,
deaths and survival, valid transfers and requests, successful synchronized
actions, and episode length.  Protocol metrics include:
\begin{itemize}
  \item plan-call validity, \code{KEEP}/\code{REPLACE} counts, coverage, age,
  expiry, version churn, and assignment switches;
  \item assignment acknowledgement rate and latency;
  \item status validity, state distribution, stale/wrong-task rejection, waiting,
  and blocker/completion/emergency-to-replan latency;
  \item unauthorized plan attempts and the stronger invariant that zero are
  accepted; and
  \item semantic compliance sampled from action traces: whether behavior advances
  the assigned task when locally feasible.
\end{itemize}

Efficiency metrics separate action calls from commander-plan calls, logical
responses from provider attempts, input/output/reasoning/cache tokens, emitted
payload bytes from fan-out delivery bytes, serial plan latency from parallel
worker-round latency, and total wall time per tick.  The sum of model-call
latencies is not substituted for wall time because calls overlap.  Failed
attempts remain in cost/error accounting while invalid episodes are excluded
from reward aggregates.

The key causal diagnostics are as important as the headline score.  If plan
validity is high but coverage is low, lease/review logic is at fault.  If coverage
is high but acknowledgement is low, instructions may not be salient.  If
acknowledgement is high but assigned behavior is absent, the result is linguistic
compliance without execution.  If achievement improves but token/wall cost breaks
the gate, the mechanism is effective but not scalable.  If star changes behavior
while plan traces remain fixed, peer status visibility---not command authority---
is the likely cause.

\section{Extension to multiple squads}

The current schema includes \code{team\_id} now, but implementing several squads
is intentionally deferred.  A scalable extension requires decisions that the
one-squad study cannot answer:
\begin{enumerate}
  \item whether squads are fixed or formed from capability evidence;
  \item whether commanders negotiate peer-to-peer, elect a mission commander,
  or use a task market;
  \item which status summaries cross squad boundaries and at what rate;
  \item how cross-squad dependencies, leases, and cancellation propagate; and
  \item what happens when a commander fails, partitions, or equivocates.
\end{enumerate}

A conservative design is a two-level namespace.  Each local commander owns only
its squad plan; an inter-squad record may request or promise a capability but
cannot directly assign another squad's wingmen.  Cross-squad work uses explicit
dependencies and leases, and raw local observations never flow automatically to
another planner.  Evaluate $N=3,6,12,24$ agents with one squad per three agents,
then compare commander mesh, sparse task-scoped routing, and a bodyless mission
planner as separate treatments.  Communication cost must be normalized by
deliveries rather than emitted messages.

\section{Threats to validity and falsification criteria}

The most serious threat is treatment bundling.  The full intervention adds a
planner call, structured plan, authority instruction, status format, and event
logic.  The first two-arm study tests the package.  Subsequent ablations are
needed to locate causality, but adding a prompt-only intermediate arm to the first
paid stage would dilute the locked question and budget.

Model stochasticity and three development seeds limit inference.  Agent 0 has a
particular specialization and spawn state.  Structured output may consume tokens
or distract from action tags.  A strict parser can undercount semantically useful
free text, while a permissive parser can inflate compliance.  Status evidence is
self-reported; validation establishes authority and freshness, not truth.  A
commander may hallucinate a coherent but infeasible plan from its partial view.
Some episodes may contain few opportunities for genuinely joint achievements.
Manual semantic review is therefore required alongside automated metrics.

The hypothesis is falsified, rather than rescued by post-hoc narrative, if the
frozen treatment cannot maintain plan/status validity, causes material action
format degradation, produces assignments that are acknowledged but not executed,
fails to improve held-out outcomes, or exceeds the preregistered overhead bound.
Positive plan-coherence metrics with no task effect would support the mechanism
claim but not the performance claim.  Positive results only on commander ID 0
would indicate specialization or position dependence, not a general hierarchy
effect.

\section{Current status and reproducibility contract}

The implementation branch contains the new state machine, runtime integration,
profiles, launcher, summarizer, tests, and this report.  New hierarchy outcomes
are deliberately marked \emph{pending}.  Each paid study root is required to
contain source and lock hashes, resolved base and per-arm configuration hashes,
the prompt-contract hash, exact commands and cache keys, seed/call ceilings,
canonical episode artifacts, append-only attempt ledgers, debug traces, a machine
summary, and a Markdown report.  A dirty launch requires an explicit override and
records the status; existing outputs are never relabeled as a new protocol.

The execution order is:
\begin{verbatim}
./commands.sh commander-study --stage 30 --dry-run
./commands.sh commander-study --stage 30 --preflight
./commands.sh commander-study --stage 30
# inspect automated gate and semantic traces
./commands.sh commander-study --stage 200 --dry-run
./commands.sh commander-study --stage 200
\end{verbatim}
The final two commands are run only after the short gate.  \code{--include-star}
adds the routing ablation on a fresh study; it does not alter the preregistered
Source-versus-broadcast contrast.

\section{Conclusion}

The best first hierarchy experiment is not a larger committee of planners.  It
is a small, auditable separation of concerns: one locally informed embodied
commander owns a short-lived whole-team plan; every agent, including the
commander, executes through the unchanged grounded action policy; wingmen retain
tactical autonomy; and validated reports trigger bounded replanning.  Classical
teamwork theory supplies the commitment vocabulary, MARL supplies temporal role
abstraction, LLM systems supply flexible semantic decomposition, embodied
planning supplies the grounding warning, and distributed leases supply a useful
stale-authority discipline.  The staged experiment is designed to discover
whether these pieces produce actual collective performance rather than merely
better-formatted agreement.

\bibliographystyle{plainnat}
\bibliography{references}
\end{document}
